\chapter{Mirror-Channel TBA and Finite-Size Effects}
\label{ch:integrable-mirror-channel-tba-finite-size-wrapping}

\ConventionsMixedCFT

This chapter refines the finite-volume interpretation of the thermodynamic
Bethe ansatz.  The primitive object is a Euclidean two-torus partition
function with two possible Hamiltonian decompositions.  The mirror-channel
description then turns finite circumference in the original theory into
finite temperature in the rotated theory.

\section{Two Hamiltonian Decompositions}

Let the Euclidean spacetime be a flat torus with cycles of lengths \(R\) and
\(L\).  Quantization along the \(L\)-direction gives
\[
  Z(R,L)=\operatorname{Tr}_{\Hilb_R}\exp[-L H_R],
\]
where \(H_R\) is the Hamiltonian on a spatial circle of circumference \(R\).
If the theory has a unique ground state on that circle and a spectral gap
above it, then
\[
  -\lim_{L\to\infty}\frac{1}{L}\log Z(R,L)=E_0(R).
\]
Quantization after exchanging the two cycles gives
\[
  Z(R,L)=\operatorname{Tr}_{\widetilde{\Hilb}_L}\exp[-R\widetilde H_L].
\]
The Hilbert space \(\widetilde{\Hilb}_L\) is the mirror-channel Hilbert space.
For a relativistic theory obtained by Euclidean rotation, the one-particle
dispersion in the infinite-volume mirror channel has the same form
\[
  \widetilde E_a(\theta)=m_a\cosh\theta,\qquad
  \widetilde p_a(\theta)=m_a\sinh\theta,
\]
but the physical interpretation of \(R\) changes: it is the inverse
temperature of the mirror gas.

\begin{proposition}[Vacuum energy from mirror pressure]
\label{prop:mirror-pressure-vacuum-energy}
Assume that the two torus decompositions above exist, the large-\(L\) limit
of \(L^{-1}\log Z(R,L)\) exists, and the mirror channel has a thermodynamic
limit at inverse temperature \(R\).  Let \(f_{\mathrm{mir}}(R)\) be the
mirror free-energy density, defined by
\[
  f_{\mathrm{mir}}(R)
  =
  -\lim_{L\to\infty}\frac{1}{RL}\log
  \operatorname{Tr}_{\widetilde{\Hilb}_L}\exp[-R\widetilde H_L].
\]
Then the direct-channel vacuum energy is
\[
  E_0(R)=R\,f_{\mathrm{mir}}(R).
\]
\end{proposition}

\begin{proof}
Both traces compute the same finite-regulator torus functional.  Taking
the logarithm and dividing by \(L\) gives
\[
  \frac{1}{L}\log Z(R,L)
  =
  \frac{1}{L}\log
  \operatorname{Tr}_{\widetilde{\Hilb}_L}\exp[-R\widetilde H_L].
\]
The direct-channel large-\(L\) limit equals \(-E_0(R)\).  The mirror
thermodynamic definition gives the limit on the right as
\(-R f_{\mathrm{mir}}(R)\).  Equating the two limits proves the formula.
\end{proof}

\section{Mirror TBA Functional}

For diagonal factorized scattering with species \(a\), masses \(m_a\), and
two-particle phase shifts
\[
  \varphi_{ab}(\theta)=-\frac{1}{2\pi\ii}\frac{d}{d\theta}
  \log S_{ab}(\theta),
\]
the mirror thermodynamic variational problem uses particle and hole rapidity
densities \(\rho_a(\theta)\) and \(\rho_a^h(\theta)\) obeying
\[
  \rho_a(\theta)+\rho_a^h(\theta)
  =
  \frac{m_aL}{2\pi}\cosh\theta
  +
  \sum_b\int_{\R}d\theta'\,
  \varphi_{ab}(\theta-\theta')\rho_b(\theta').
\]
The entropy density of a fermionic Bethe cell is
\[
  s_a=
  (\rho_a+\rho_a^h)\log(\rho_a+\rho_a^h)
  -\rho_a\log\rho_a-\rho_a^h\log\rho_a^h .
\]
Extremizing mirror free energy at inverse temperature \(R\) gives
\[
  \epsilon_a(\theta)
  =
  Rm_a\cosh\theta
  -
  \sum_b\int_{\R}d\theta'\,
  \varphi_{ab}(\theta-\theta')
  \log\!\left(1+e^{-\epsilon_b(\theta')}\right),
\]
where \(e^{\epsilon_a}=\rho_a^h/\rho_a\).  Substitution into
Proposition~\ref{prop:mirror-pressure-vacuum-energy} gives
\[
  E_0(R)=\mathcal E_{\mathrm{bulk}}R
  -
  \sum_a\int_{\R}\frac{d\theta}{2\pi}\,
  m_a\cosh\theta\,
  \log\!\left(1+e^{-\epsilon_a(\theta)}\right).
\]
The constant \(\mathcal E_{\mathrm{bulk}}\) is the local extensive coordinate
already encountered in Chapter~\ref{ch:integrable-rg-flows-perturbed-cft}.

\section{Large-Circumference Expansion}

Let \(m_{\min}\) be the smallest particle mass.  If \(Rm_{\min}\gg1\), then
the TBA equation may be solved by iteration in the functions
\[
  q_a(\theta)=\exp[-Rm_a\cosh\theta].
\]
The first term is
\[
  \epsilon_a^{(0)}(\theta)=Rm_a\cosh\theta .
\]
Inserting this into the energy functional gives the single-winding vacuum
correction
\[
  E_0(R)-\mathcal E_{\mathrm{bulk}}R
  =
  -\sum_a\int_{\R}\frac{d\theta}{2\pi}\,
  m_a\cosh\theta\,q_a(\theta)
  +O(q^2),
\]
where \(O(q^2)\) denotes terms containing at least two mirror occupation
factors.  This statement is an asymptotic expansion at fixed scattering
kernel and fixed mass spectrum.  A quantitative estimate requires a norm on
the kernels and a lower bound on \(Rm_{\min}\).

\section{Excited States}

A finite-volume excited state in the direct channel is described by rapidities
\(\theta_j\) and species \(a_j\) satisfying Bethe-Yang equations corrected by
virtual mirror particles.  In a diagonal theory without bound-state
singularities crossing the integration contour, the first mirror correction
has the form
\[
  \Delta E_{\Psi}^{F}(R)
  =
  -\sum_b\int_{\R}\frac{d\theta}{2\pi}\,
  m_b\cosh\theta\,e^{-Rm_b\cosh\theta}
  \left(
  \prod_j S_{b a_j}\!\left(\theta-\theta_j+\frac{\ii\pi}{2}\right)
  -1
  \right),
\]
with all rapidities and scattering phases understood in the analytic strip
where the displayed contour is valid.  The term \(-1\) subtracts the vacuum
mirror loop.  Bound-state poles generate additional residue terms when the
contour crosses them; their coefficients are fixed by the same three-point
couplings that enter the factorized scattering bootstrap.

\section{Wrapping Effects}

A wrapping contribution is a term in a finite-volume observable represented
in the mirror channel by virtual propagation around a compact cycle.  The
order is measured by the number and species of mirror occupation factors
\(q_a(\theta)\), not by the loop number of a Lagrangian perturbation series.
For a local operator \(\widehat{\mathcal O}\), a finite-volume diagonal matrix
element has an expansion
\[
  \langle \Psi,R|\widehat{\mathcal O}|\Psi,R\rangle
  =
  \langle \Psi|\widehat{\mathcal O}|\Psi\rangle_{\infty}
  +
  \sum_{k\geq1}\Delta_{\mathcal O}^{(k)}(R),
\]
where \(\Delta_{\mathcal O}^{(k)}\) is expressed through \(k\)-mirror-particle
form factors and finite-volume density factors when the form-factor series
converges in the required strip.  This formula is useful only after the
operator normalization and contour prescription have been declared.

\section{Relation To Planar Wrapping}

In large-\(N\) gauge-theory spectral problems the word wrapping often refers
to corrections whose range reaches the length of a single-trace operator.
The relativistic finite-size mechanism developed here is a different
construction: the compact cycle is an actual spacetime circle of a
two-dimensional QFT, and the mirror particles are states in the rotated
finite-temperature channel.  Later planar-integrability chapters must build
an explicit bridge between these two settings by identifying the spectral
problem, its worldsheet or spin-chain mirror channel, and the analytic
continuation that maps finite-size corrections between them.

\begin{openproblem}[Excited-state mirror construction]
\label{op:excited-state-mirror-construction}
Give a theorem-level construction of excited-state TBA starting from
factorized scattering, a declared analytic strip, pole data, and finite-volume
Bethe states, including a proof that the contour deformations and residue
terms reproduce the spectral trace on the torus.
\end{openproblem}
